\section{Move Semantics}

Move semantics, introduced in C++11, is one of the most significant additions to modern C++. It enables efficient transfer of resources from one object to another, avoiding expensive copies when they're not needed.

\subsection{The Problem: Expensive Copies}

Consider returning a large vector from a function:

\begin{lstlisting}
	std::vector<int> createLargeVector() {
		std::vector<int> result(1000000);  // 1 million integers
		// ... fill the vector ...
		return result;  // Pre-C++11: potentially copies 1 million integers!
	}
	
	std::vector<int> v = createLargeVector();
\end{lstlisting}

Before C++11, returning \texttt{result} could involve copying all one million integers to a new location. Move semantics allows us to simply "steal" the internal resources instead.

\subsection{Lvalues and Rvalues}

To understand move semantics, you must first understand the difference between \textbf{lvalues} and \textbf{rvalues}.

\subsubsection{Lvalues}

An \textbf{lvalue} (locator value) is an expression that refers to a memory location. You can take its address with \texttt{\&}.

\begin{lstlisting}
	int x = 10;         // x is an lvalue
	int* p = &x;        // We can take its address
	
	std::string s = "hello";  // s is an lvalue
\end{lstlisting}

\subsubsection{Rvalues}

An \textbf{rvalue} is a temporary value that doesn't persist beyond the expression that uses it. You cannot take its address.

\begin{lstlisting}
	int x = 10 + 20;    // (10 + 20) is an rvalue - temporary result
	// int* p = &(10 + 20);  // ERROR: can't take address of rvalue
	
	std::string getName() { return "Alice"; }
	std::string s = getName();  // The returned string is an rvalue
\end{lstlisting}

\subsubsection{Why This Matters}

Rvalues are temporaries that are about to be destroyed. If we're copying from an rvalue, we can safely "steal" its resources instead of copying—the temporary won't need them anymore.

\subsection{Rvalue References}

C++11 introduced \textbf{rvalue references}, declared with \texttt{\&\&}:

\begin{lstlisting}
	int x = 10;
	int& lref = x;       // Lvalue reference: binds to lvalues
	int&& rref = 20;     // Rvalue reference: binds to rvalues
	
	// int& lref2 = 20;     // ERROR: lvalue ref can't bind to rvalue
	// int&& rref2 = x;     // ERROR: rvalue ref can't bind to lvalue
\end{lstlisting}

Rvalue references allow us to overload functions to behave differently for temporaries.

\subsection{Move Constructor and Move Assignment}

\subsubsection{The Move Constructor}

A \textbf{move constructor} transfers resources from a temporary (rvalue) instead of copying:

\begin{lstlisting}
	class String {
		private:
		char* data;
		size_t length;
		
		public:
		// Copy constructor (deep copy)
		String(const String& other) 
		: length(other.length), data(new char[length + 1]) {
			std::strcpy(data, other.data);
			std::cout << "Copy constructor called\n";
		}
		
		// Move constructor (steal resources)
		String(String&& other) noexcept
		: data(other.data), length(other.length) {
			other.data = nullptr;   // Leave source in valid state
			other.length = 0;
			std::cout << "Move constructor called\n";
		}
	};
\end{lstlisting}

\subsubsection{The Move Assignment Operator}

\begin{lstlisting}
	class String {
		public:
		// Move assignment operator
		String& operator=(String&& other) noexcept {
			if (this != &other) {
				delete[] data;          // Release current resources
				
				data = other.data;      // Steal resources
				length = other.length;
				
				other.data = nullptr;   // Leave source in valid state
				other.length = 0;
			}
			std::cout << "Move assignment called\n";
			return *this;
		}
	};
\end{lstlisting}

\subsubsection{When Move Operations Are Called}

\begin{lstlisting}
	String createString() {
		return String("Hello");  // Returns an rvalue
	}
	
	int main() {
		String s1("Hello");           // Normal constructor
		String s2 = s1;               // Copy constructor (s1 is lvalue)
		String s3 = createString();   // Move constructor (rvalue returned)
		
		String s4("World");
		s4 = s1;                      // Copy assignment (s1 is lvalue)
		s4 = createString();          // Move assignment (rvalue returned)
	}
\end{lstlisting}

\subsection{std::move}

\texttt{std::move} is a cast that converts an lvalue to an rvalue reference, enabling move operations:

\begin{lstlisting}
	#include <utility>  // for std::move
	
	String s1("Hello");
	String s2 = s1;              // Copy: s1 is lvalue
	String s3 = std::move(s1);   // Move: std::move casts s1 to rvalue
	
	// WARNING: s1 is now in a "moved-from" state!
	// It's valid but unspecified - don't use its value
\end{lstlisting}

\begin{remark}
	\texttt{std::move} doesn't actually move anything—it's just a cast. The actual move happens when the move constructor or move assignment operator is called. After using \texttt{std::move} on an object, treat it as if it no longer has a useful value.
\end{remark}

\subsection{The Rule of Five}

If you define any of these special member functions, you should consider defining all five:

\begin{enumerate}
	\item Destructor
	\item Copy constructor
	\item Copy assignment operator
	\item Move constructor
	\item Move assignment operator
\end{enumerate}

\begin{lstlisting}
	class Resource {
		private:
		int* data;
		
		public:
		// Constructor
		Resource(int value) : data(new int(value)) {}
		
		// 1. Destructor
		~Resource() { delete data; }
		
		// 2. Copy constructor
		Resource(const Resource& other) : data(new int(*other.data)) {}
		
		// 3. Copy assignment
		Resource& operator=(const Resource& other) {
			if (this != &other) {
				delete data;
				data = new int(*other.data);
			}
			return *this;
		}
		
		// 4. Move constructor
		Resource(Resource&& other) noexcept : data(other.data) {
			other.data = nullptr;
		}
		
		// 5. Move assignment
		Resource& operator=(Resource&& other) noexcept {
			if (this != &other) {
				delete data;
				data = other.data;
				other.data = nullptr;
			}
			return *this;
		}
	};
\end{lstlisting}

\subsection{The Rule of Zero}

Better yet, if you use RAII types (smart pointers, standard containers), you often don't need to define any special member functions:

\begin{lstlisting}
	class Person {
		private:
		std::string name;                    // Handles its own memory
		std::vector<std::string> friends;    // Handles its own memory
		std::unique_ptr<Address> address;    // Handles its own memory
		
		public:
		Person(std::string n) : name(std::move(n)) {}
		
		// No destructor needed
		// No copy/move constructors needed
		// No copy/move assignment needed
		// The compiler-generated versions do the right thing!
	};
\end{lstlisting}

\begin{remark}
	The Rule of Zero should be your default approach. Only implement custom special member functions when you're directly managing resources (raw pointers, file handles, etc.), and even then, consider wrapping them in RAII classes.
\end{remark}

\subsection{Perfect Forwarding}

\textbf{Perfect forwarding} preserves the value category (lvalue/rvalue) when passing arguments through template functions:

\begin{lstlisting}
	template <typename T>
	void wrapper(T&& arg) {
		// std::forward preserves whether arg was lvalue or rvalue
		actualFunction(std::forward<T>(arg));
	}
	
	void actualFunction(const std::string& s);  // For lvalues
	void actualFunction(std::string&& s);       // For rvalues
	
	std::string s = "hello";
	wrapper(s);              // Calls lvalue version
	wrapper(std::string());  // Calls rvalue version
\end{lstlisting}

\subsection{Best Practices}

\begin{enumerate}
	\item \textbf{Mark move operations \texttt{noexcept}}: This enables important optimizations, especially with standard containers
	
	\item \textbf{Leave moved-from objects in a valid state}: They should be destructible and assignable, even if their value is unspecified
	
	\item \textbf{Use \texttt{std::move} for the last use of a local variable}: When passing to a function or returning
	
	\item \textbf{Don't \texttt{std::move} return values}: The compiler handles this automatically (RVO/NRVO)
	
	\item \textbf{Prefer the Rule of Zero}: Use smart pointers and standard containers to avoid manual resource management
\end{enumerate}

\begin{lstlisting}
	// GOOD: Move on last use
	void processData(std::vector<int> data);
	
	std::vector<int> v = {1, 2, 3};
	// ... use v ...
	processData(std::move(v));  // Last use of v
	
	// BAD: Don't move return values
	std::vector<int> createVector() {
		std::vector<int> v = {1, 2, 3};
		return std::move(v);  // WRONG: prevents RVO optimization
		return v;             // RIGHT: compiler optimizes this
	}
\end{lstlisting}